%% ma: L12-B09-0002
%% lop: 12
%% chuong: 3
%% bai: 9
%% dang: 12.9.2
%% loai: TN4
%% muc: TH
%% dap_an: "D"
%% nguon: "(NBV)-ĐỀ SỐ 4-MỖI NGÀY 1 ĐỀ THI 2025.pdf (D04, MỖI NGÀY 1 ĐỀ - Đề số 4), Phần I Câu 5"
%% kien_thuc: "Bài 9 (khoảng biến thiên của mẫu số liệu ghép nhóm)"
%% trang_thai: cho-duyet
%% ly_do: "Dạng đề xuất mới 12.9.2 chờ giáo viên duyệt"
%% ghi_chu: "Bẫy: nếu không bỏ các nhóm đầu/cuối có tần số 0 thì R_A=R_B=1 và chọn nhầm A. Đáp án gốc D đúng."
\begin{ex}
Cho hai mẫu số liệu ghép nhóm $A$ và $B$ có bảng tần số ghép nhóm như sau:
\begin{center}
\renewcommand{\arraystretch}{1.25}
\begin{tabular}{|c|c|c|c|c|c|}
\hline
$A$: Nhóm & $[2{,}8;3{,}0)$ & $[3{,}0;3{,}2)$ & $[3{,}2;3{,}4)$ & $[3{,}4;3{,}6)$ & $[3{,}6;3{,}8)$\\
\hline
Tần số & $2$ & $5$ & $8$ & $5$ & $0$\\
\hline
\end{tabular}

\medskip
\begin{tabular}{|c|c|c|c|c|c|}
\hline
$B$: Nhóm & $[2{,}8;3{,}0)$ & $[3{,}0;3{,}2)$ & $[3{,}2;3{,}4)$ & $[3{,}4;3{,}6)$ & $[3{,}6;3{,}8)$\\
\hline
Tần số & $0$ & $0$ & $6$ & $4$ & $2$\\
\hline
\end{tabular}
\end{center}
Gọi $R_A$ và $R_B$ lần lượt là khoảng biến thiên của mẫu số liệu ghép nhóm $A$ và $B$. Phát biểu nào sau đây đúng?
\choice{$\dfrac{R_A}{R_B}=1$}{$\dfrac{R_A}{R_B}=\dfrac32$}{$\dfrac{R_A}{R_B}=\dfrac34$}{\True $R_A-R_B=0{,}2$}
\loigiai{Khoảng biến thiên là hiệu giữa đầu mút phải của nhóm cuối cùng có chứa dữ liệu và đầu mút trái của nhóm đầu tiên có chứa dữ liệu.
Mẫu $A$: nhóm đầu $[2{,}8;3{,}0)$, nhóm cuối có dữ liệu $[3{,}4;3{,}6)$ (nhóm $[3{,}6;3{,}8)$ có tần số $0$) nên $R_A=3{,}6-2{,}8=0{,}8$.
Mẫu $B$: nhóm đầu có dữ liệu $[3{,}2;3{,}4)$, nhóm cuối $[3{,}6;3{,}8)$ nên $R_B=3{,}8-3{,}2=0{,}6$.
Vậy $R_A-R_B=0{,}2$ và $\dfrac{R_A}{R_B}=\dfrac43$, nên chỉ D đúng.}
\end{ex}
